4장의 분석은 컴파일러 패스가 넘을 수 없는 간극이 실재함을 보인다. 미정의 동작 뒤에 숨은
하드웨어 의미론(시프트 포화), 대상 코어의 사이클 모델, 레지스터 압박과 언롤의 상호작용은
각각 프런트엔드, 백엔드, 할당기에 흩어져 있어 단일 패스가 종합적으로 추론하기 어렵다.

본 연구는 이 간극을 LLM으로 메우는 가능성을 실증한다. 실제로 본 논문의 sum\_dual\_asm은
LLM이 주도하는 반복 피드백 루프---실기 측정, 병목 분석, 코드 수정, 재측정---를 통해
도출되었다. 이 과정에서 LLM은 (1) ISA 문서의 포화 의미론과 C 표준의 미정의 동작 규칙을
대조하여 ``컴파일러가 규칙상 생성할 수 없는 최적 코드''의 존재를 예측하였고, (2) M0+
사이클 모델로 이론 성능(10.25사이클)을 사전 계산해 실측(10.0)과 2\% 이내로 일치시켰으며,
(3) 어셈블리 조각 방식의 실패(레지스터 스필)를 분석해 루프 단위 재작성으로 전환하였다.
정합성은 매 반복 하드웨어에서 기준 구현과의 결과 비교로 검증되므로, LLM의 생성 오류는
루프 안에서 걸러진다.

이 접근이 성립하는 조건도 실험에서 드러난다. LLM이 제안한 첫 어셈블리는 어셈블러 문법
불일치로 두 차례 컴파일에 실패했고, 어셈블리 조각 방식은 성능 목표에 미달했다. 그럼에도
최종 구현에 도달할 수 있었던 것은 컴파일 오류, 실측 성능, 정합성 검증이라는 세 겹의
객관적 피드백이 매 반복 제공되었기 때문이다. 즉 LLM 기반 최적화의 신뢰성은 모델 단독의
정확성이 아니라 하드웨어가 포함된 폐루프 구조에서 나온다.

이는 두 방향으로 확장될 수 있다. 첫째, LLM을 \textbf{패스 한계 예측기}로 사용하여, 주어진
커널과 대상 코어에 대해 ``컴파일러 출력과 이론 하한의 격차''를 자동 추정하고 수동 최적화가
필요한 지점만 선별하는 것이다. 둘째, LLM이 발견한 변환---예컨대 부호 가변 시프트를 식
(\ref{eq:dual})의 포화 이중 시프트로 바꾸는 것---을 페리홀(peephole) 패스나 명령어 선택
패턴으로 형식화하여 컴파일러에 환류하는 것이다. 전자는 개발자 노력의 배분을, 후자는 발견된
지식의 재사용을 자동화한다. 하드웨어 실측이 폐루프에 포함되어 있다는 점에서, 이 접근은
정확성 보장이 어려운 LLM 단독 코드 생성과 구별된다.
